笛卡尔坐标直接表示原子在空间中的位置，建模界面通常以埃（\AA）为单位。分数坐标则表示沿三条晶格矢量的比例。若分数坐标为 $(u,v,w)$，原子的位置就是 $u\mathbf a+v\mathbf b+w\mathbf c$。这三个数没有长度单位；它们回答的是“沿每条晶格矢量走多少”，而不是“沿全局坐标轴走多少埃”。

为读懂晶格输入表格，可以把三条矢量按行排成一个矩阵：
\[
 A=\begin{bmatrix}
 a_x&a_y&a_z\\
 b_x&b_y&b_z\\
 c_x&c_y&c_z
 \end{bmatrix},\qquad
 \mathbf f=(u,v,w),\qquad \mathbf r=(x,y,z).
\]
矩阵的第一行是 $\mathbf a$，第二、三行依次是 $\mathbf b$、$\mathbf c$；每行中的三列才是全局 $x,y,z$ 分量。本书将坐标写成行向量，因此两种坐标的转换为
\begin{equation}
 \mathbf r=\mathbf f A,\qquad \mathbf f=\mathbf r A^{-1}.
 \label{eq:row-lattice-coordinates}
\end{equation}
这里 $A^{-1}$ 是整个晶格矩阵的逆矩阵，不能把表格中的每个数分别取倒数。晶格三行所定义的晶胞应有非零体积，才能由笛卡尔位置求回唯一的分数坐标。

以一组便于手算的晶格说明：令 $\mathbf a=(4,0,0)$、$\mathbf b=(1,3,0)$、$\mathbf c=(0,0,6)$，单位均为\AA。分数坐标 $(0.5,0.5,0)$ 对应
\[
 \mathbf r=0.5(4,0,0)+0.5(1,3,0)
          =(2.5,1.5,0)\,\text{\AA}.
\]
沿 $\mathbf b$ 走一半，同时包含全局 $x$ 方向的 $0.5\,\text{\AA}$ 和 $y$ 方向的 $1.5\,\text{\AA}$。所以斜晶胞中的第二个分数分量，不能直接当作全局 $y$ 坐标。图\ref{fig:round10-fractional-cartesian}将这一次相加画成两段相接的箭头。

\begin{figure}[!htb]
\centering
\begin{tikzpicture}[x=1mm,y=1mm,font=\sffamily\small]
 \fill[paleblue] (8,6)--(56,6)--(68,42)--(20,42)--cycle;
 \draw[teal,line width=.7pt] (8,6)--(56,6)--(68,42)--(20,42)--cycle;
 \draw[->,>=stealth,muted] (8,6)--(84,6) node[right] {$x$};
 \draw[->,>=stealth,muted] (8,6)--(8,49) node[above] {$y$};
 \draw[->,>=stealth,brand,line width=1pt] (8,6)--(32,6);
 \draw[->,>=stealth,brand,line width=1pt] (32,6)--(38,24);
 \draw[->,>=stealth,teal,line width=.8pt] (8,6)--(38,24);
 \draw[dashed,muted] (38,24)--(38,6);
 \fill[brand] (38,24) circle (1.2);
 \node[anchor=north,text=brand] at (22,5) {$u\mathbf a$};
 \node[anchor=west,text=brand] at (38,14) {$v\mathbf b$};
 \node[anchor=south east,text=teal] at (25,16) {$\mathbf r$};
 \node[anchor=north east] at (7,5) {$O$};
 \node[anchor=west,text=brand] at (42,27) {$\mathbf f=(u,v,0)$};
 \node[anchor=west,text=muted,font=\sffamily\footnotesize] at (78,35) {晶胞在平面内的截面};
 \node[anchor=west,text=muted,font=\sffamily\footnotesize] at (78,28) {箭头相加得到位置};
 \node[anchor=north,text=teal] at (54,4) {$\mathbf a$};
 \node[anchor=east,text=teal] at (15,31) {$\mathbf b$};
\end{tikzpicture}
\caption{分数坐标与笛卡尔位置的关系（二维示意）。$u\mathbf a$ 与 $v\mathbf b$ 相加得到 $\mathbf r$；斜晶胞的 $\mathbf b$ 方向同时含有全局 $x$ 和 $y$ 分量。}
\label{fig:round10-fractional-cartesian}
\end{figure}


在边长为 $5\,\text{\AA}$ 的立方晶胞中，$(0.5,0,0)$ 则对应沿 $\mathbf a$ 前进 $2.5\,\text{\AA}$。两种例子的区别来自晶格，而不是坐标转换规则不同。编辑时若保持 $\mathbf f$ 不变而改变 $A$，原子通常会随晶胞移动；若希望原子在空间中的位置不动，则需保留 $\mathbf r$，并用新晶格重新求 $\mathbf f$。第\ref{ch:modeling-tools}章的\ref{sec:lattice-transform-reference}节介绍这两个选项。

\begin{table}[!htb]
\centering\small
\begin{tabularx}{\linewidth}{llX}
\toprule
符号 & 单位 & 在参数中的含义\\
\midrule
$A$ & \AA & 晶格矩阵，三行依次为 $\mathbf a,\mathbf b,\mathbf c$。\\
$\mathbf f$ & 无量纲 & 沿三条晶格矢量的比例，即分数坐标。\\
$\mathbf r,\mathbf p$ & \AA & 笛卡尔位置，以及旋转锚点或镜面上的点。\\
$\mathbf n,\hat{\mathbf n}$ & 方向 & 非零方向向量及其单位向量；只看方向时不以\AA 填长度。\\
$\theta$ & 度 & 对话框中的转角，正角按右手定则。\\
$S$ & 整数 & 定义新晶格的超胞变换矩阵。\\
$N,V$ & 个、\AA$^3$ & 模型独立位点数与晶胞体积。\\
$\varepsilon_a,\varepsilon_b,\varepsilon_c$ & 无量纲 & 沿三条晶格矢量的相对伸缩量。\\
\bottomrule
\end{tabularx}
\caption{阅读建模公式时常用的符号。坐标和锚点使用长度，方向向量只指定朝向；两者在对话框中有不同作用。}
\label{tab:round10-modeling-symbols}
\end{table}
